\section{Introduction}\label{sec:introduction}

For a finite set $P\subset\R^2$ and a point $x\in P$, write
\[
D_x(P)=\{\lVert y-x\rVert_2:y\in P\setminus\{x\}\}.
\]
The weak pinned Erd\H{o}s distinct-distance conjecture asserts that, for
every $\varepsilon>0$, every sufficiently large $n$-point planar set has
a point $x$ with $|D_x(P)|\ge n^{1-\varepsilon}$. Erd\H{o}s explicitly
stated this question in 1957 \cite[Problem~16]{Erdos1957}; the modern
terminology appears, for example, in
\cite[Conjecture~4.7]{WeakPinned2025}. We resolve it positively by
proving a uniform assertion about distance multiplicities at their
individual pins. Its consequence will give the stated number of
distances at all but $o(n)$ pins.

For distinct $x,y\in P$, define the \emph{distance fiber size}
\[
k_P(x,y)=\bigl|\{z\in P\setminus\{x\}:
                  \lVert z-x\rVert_2=\lVert y-x\rVert_2\}\bigr|.
\]
This counts $y$ itself. Different pairs use their own source point $x$;
the definition does not pool pairs having one common numerical distance.
For $n\ge2$ and $s>0$, put
\[
F_n(s)=\sup_{\substack{P\subset\R^2\\|P|=n}}
 \frac{\bigl|\{(x,y)\in P^2:x\ne y,\ k_P(x,y)\ge n^s\}\bigr|}{n(n-1)}.
\]

\begin{theorem}\label{thm:main}
For every fixed $s>0$, one has $F_n(s)\longrightarrow0$ as
$n\longrightarrow\infty$.
\end{theorem}

No hypothesis on separation, general position, or the coordinates is
imposed. The convergence is uniform over all configurations. Its
exponent $s$ is fixed; the theorem supplies no explicit rate of decay.

\begin{corollary}\label{cor:pins}
For every fixed $\varepsilon>0$,
\[
\sup_{\substack{P\subset\R^2\\|P|=n}}
 \frac{\bigl|\{x\in P:|D_x(P)|<n^{1-\varepsilon}\}\bigr|}{n}
 \longrightarrow0.
\]
In particular, every sufficiently large $n$-point planar set has a pin
determining at least $n^{1-\varepsilon}$ distinct nonzero distances.
\end{corollary}

\begin{proof}
For $\varepsilon\ge1$, every pin in a set of at least two points already
has at least one distance. Suppose $0<\varepsilon<1$, and take
$s=\varepsilon/2$. At a pin with fewer than $n^{1-\varepsilon}$ distance
classes, fewer than $n^{1-\varepsilon+s}$ neighbors lie in classes of
size less than $n^s$. Each such pin therefore contributes at least
$n-1-n^{1-\varepsilon/2}$ pairs counted by $F_n(s)$. If there are $b$
such pins, then, for all sufficiently large $n$,
\[
\frac bn\le
 \frac{(n-1)F_n(\varepsilon/2)}{n-1-n^{1-\varepsilon/2}}
 \longrightarrow0.
\]
The bound is independent of $P$.
\end{proof}

\subsection{Background and significance}
Erd\H{o}s initiated the planar distinct-distance and repeated-distance
problems in 1946 \cite{Erdos1946}. The global distinct-distance problem
asks for the minimum size of the union $\bigcup_{x\in P}D_x(P)$.
The square lattice gives sets for which this union has size at most a
constant multiple of $n/\sqrt{\log n}$ \cite[Theorem~1]{Erdos1946}.
Guth and Katz proved that every planar $n$-point
set determines at least a constant multiple of $n/\log n$ distances
\cite[Theorem~1.1]{GuthKatz}. Their conclusion concerns the union;
the pinned question requires the distances to have a common source.

The development of pinned bounds combined incidence geometry with
increasingly effective counting and entropy estimates. Solymosi and
T\'oth proved a pinned lower bound of order $n^{6/7}$
\cite[Theorem~1]{SolymosiToth}. Tardos improved this exponent using
entropy inequalities for patterns of sums and differences
\cite{Tardos2003}. Katz and Tardos refined that method to obtain every
exponent below
\[
 \frac{48-14e}{55-16e}=0.864137\ldots
\]
\cite{KatzTardos}. These results guarantee one pin in every
configuration. The sharper pinned conjecture asks for order
$n/\sqrt{\log n}$ distances at a pin; see
\cite[Introduction]{LundPetridis}. Lund and Petridis also explain the
distinction between the real planar problem and pinned-distance
theorems over other fields.

Theorem~\ref{thm:main} controls the multiplicity of a distance chosen
by sampling a distinct ordered pair. Corollary~\ref{cor:pins} makes
its almost-every-pin consequence explicit and resolves the weak
pinned conjecture. An exceptional set is necessary: if $P$ consists of
points on a circle and its center, the center sees just one nonzero distance.

\subsection{The proof in outline}
Assuming that Theorem~\ref{thm:main} fails, we first choose a sequence
with nearly maximal density of pairs in large distance fibers.
Section~\ref{sec:extremal} extracts a directed graph of these pairs
while controlling how much of a fiber can lie in a small subset.
Transfer for real closed fields preserves every distance equality and
disequality and puts each finite configuration in a number field
\cite[Theorems~3.1 and~4.1]{RCFTransfer}. Neither the field nor its
degree is fixed.

The proof then studies distance multiplicities through two coordinate maps.
For $x=(u,v)$, set $Z_1(x)=u+iv$ and $Z_2(x)=u-iv$. Squared distance
factors as
\[
\lVert x-y\rVert_2^2
   =(Z_1(y)-Z_1(x))(Z_2(y)-Z_2(x)).
\]
Thus on a distance fiber at $x$, the second coordinate is a fractional
linear function of the first. The normalized product formula balances
the logarithms of nonzero coordinate differences across all absolute
values of the number field \cite[Theorem~7.15]{ProductFormula}.
At other complex embeddings the two maps need not be conjugates;
the algebraic factorization remains valid.

We turn this balance into an identity for finite nested partitions.
At each absolute value, points are grouped according to closeness.
Randomly shifted and rotated square grids provide nested partitions
at the complex embeddings. At any level, a cell containing more than
half the points, called a \emph{giant}, is replaced by its complement,
with multiplicities retained. The resulting overlap of two
\emph{distinct} points, summed
over all levels and absolute values, is the sum of a function of the
first point and a function of the second. This additive identity is
the main use of arithmetic in the proof.

Randomly shifted square dissections appear in Arora's Euclidean
approximation algorithms \cite[Section~2.1]{Arora1998}.
Nested partitions and their associated trees also occur in probabilistic
metric approximation, as in the work of Fakcharoenphol, Rao and Talwar
\cite{FakcharoenpholRaoTalwar2003}. Their approximation controls
multiplicative distortion. Here the required grid estimate has a
pair-independent mean additive error in logarithmic depth; this
specific property permits the product-formula cancellation.

We measure the total scale by the integrated distinct-pair overlaps
and the average lengths of levels spent outside giant cells.
The next step is to compare the uniform distribution on a fiber with
the uniform distribution on the whole configuration. The extremal
subset estimate controls cells that contain a source and a substantial
fraction of its fiber. For an actual partition cell containing only a
small fraction, the other coordinate supplies the missing control: without grid
errors, its source cell at the complementary level contains the rest
of the fiber. Together these estimates make the small-cell contribution
negligible relative to one plus the total scale
(Section~\ref{sec:smallcells}).

We then compare three laws on distinct ordered pairs: the uniform law
within a fiber, a law with one fiber-uniform and one ambient-uniform
marginal, and the uniform law on the whole configuration. The first
minus twice the second plus the third has zero total mass and zero
sum of its marginals. Its integral against the additive overlap kernel
therefore vanishes. Expanding the cell probabilities gives a squared
discrepancy between fiber mass and uniform mass, together with explicit
finite-population corrections. Large support sizes and the small-cell
estimate control these corrections. Section~\ref{sec:variance}
establishes the resulting variance bound, both for fibers and for
the sources of edges ending at a target. The use of distinct pairs is
essential: singleton cells persist through an infinite range of levels,
over which even a small diagonal error could not be integrated.

After passing to a subsequence, this nonnegative scale either grows
without bound or stays bounded. In the first case, we retain cells whose
mass and complementary mass are both bounded below. They give weighted finite rooted
trees with total length comparable to the scale.
The fiber relation and the fiber variance bound make the source-to-target distance in one tree
close to the root-to-target distance in the other. The latter depends
only on the target. Transport on the trees, using the incoming-law
variance bound, lets us replace two sources arriving at the same target
by independent uniform points. After replacing the target law as well,
the expected absolute difference of their distances to that target
would be negligible relative to the total overlap scale. An elementary
geometric bound for trees contradicts this conclusion
(Section~\ref{sec:unbounded}).

If the overlap scale stays bounded, one complex embedding can be
normalized so that the two coordinate distributions have non-atomic
limits. The fiber and incoming-neighbor estimates yield a family of
M\"obius maps carrying one limiting coordinate law to the other. The
first source coordinate is independent of the entire target pair;
the second source coordinate need not be. To remove the latter
dependence, fix a target pair and sample a point from the first limiting
coordinate law. Its image under the fiber map has the second limiting
law. The fiber equation writes the negative logarithm of the distance
from that image to the second target coordinate as a difference of
logarithmic distances from the sampled point to the first source and
target coordinates, up to an additive constant. Independence lets the
first source coordinate approach the first target coordinate. This
difference then collapses in probability under the fixed sampling law,
whereas the image's logarithmic target-distance law has positive spread
even after additive constants are discarded.
Section~\ref{sec:bounded} makes this contradiction precise.

The tree-transport formula is a standard sum of the imbalances across
tree edges \cite[Section~2, Equation~(5)]{TreeTransport}; we prove its
finite form and the conditional comparisons needed here. Weak
compactness and Portmanteau provide subsequential probability laws
\cite[Proposition~5.3 and Theorems~4.2,~3.2]{WeakCompactness}.
The local arguments must additionally establish non-atomicity,
independence of the source parameter, and preservation of the fiber
map. These properties connect the standard tools to the additive
overlap identity and its exact variance consequence.

Section~\ref{sec:extremal} prepares the extremal graph and its sampling
laws. Section~\ref{sec:hierarchies} constructs the partitions and proves
the additive identity. Sections~\ref{sec:smallcells} and
\ref{sec:variance} establish the overlap and variance estimates.
Sections~\ref{sec:unbounded} and \ref{sec:bounded} exclude the two scale
regimes and complete the proof of Theorem~\ref{thm:main}.
